\documentclass[11pt,letterpaper]{article}
\usepackage[utf8]{inputenc}
\usepackage[T1]{fontenc}
\usepackage{geometry}
\geometry{a4paper, margin=2cm}
\usepackage{hyperref}
\usepackage{setspace}
\usepackage{siunitx}
\sisetup{separate-uncertainty=true, range-phrase=--}
\DeclareSIUnit\kcalmol{kcal\,mol^{-1}}
\DeclareSIUnit\angstrom{\text{\AA}}
\singlespacing

\begin{document}
\emergencystretch=2em

\begin{flushleft}
\textbf{Myke-Vital Temgoua Sao} \\
Department of Physics, Faculty of Science \\
University of Yaound\'{e} I \\
P.O. Box 812, Yaound\'{e}, Cameroon \\
\href{mailto:myke-vital.sao@facsciences-uy1.cm}{myke-vital.sao@facsciences-uy1.cm} \\[1.5em]
August 30, 2026
\end{flushleft}

\begin{flushleft}
Prof.\ Kenneth M.\ Merz, Jr., Editor-in-Chief \\
\textit{Journal of Chemical Information and Modeling} \\
American Chemical Society
\end{flushleft}

\vspace{0.5em}

\noindent Dear Prof.\ Merz,

We submit our manuscript \textbf{``Resistance-Aware Docking Prioritization and Short MD Stress Tests Measure Non-Equivalent Quantities in an Antimalarial Candidate Cohort''} for consideration as an Article in \textit{JCIM}.

The central question is whether a resistance-aware, target-level docking workflow (RRS + PNS + ACSI) produces a coherent prioritization signal inside a curated antimalarial cohort, and whether a short single-replicate MD structural-stress test is concordant with that signal or instead reveals non-equivalence. Seventeen Set-C candidates were ranked by a docking-derived resistance-retention score across six PfDHFR and PfCRT resistance states; chemical-space (ACSI) and network (PNS) scores were kept as separate ranking dimensions. On the 12 candidates with complete two-target data the classes were one A*, one A, four B, five C, and one D. Four parent-study complexes outside the Set-C list were simulated for 10 ns each as a pose-retention check; only PfCRT--214 yielded an interpretable MM-GBSA estimate ($-18.25\pm0.40$ kcal mol$^{-1}$). A 16-system single-replicate pilot on two Set-C candidates did not recover the docking retention pattern (directional disagreement in seven of eight mutant comparisons that possessed a wild-type docking reference). An external 39-ligand panel averaged near RRS 100. The ranking is informative inside the computational panel only as a hypothesis-generating prioritization, not as a validated triage or filtering procedure; the short MD stress test demonstrates that the two quantities are non-equivalent. Neither result claims biological target engagement, affinity, or resistance circumvention.

The MD component is reported in accordance with the JCIM guidelines of Soares et al.: single-replicate limitation, force-field versions, ensemble, trajectory acceptance criteria, and data-availability statements are explicit. Companion Paper 1 supplies only upstream chemical-space context, not independent replication.

No external funding was received for this work. Code, source data, and analysis scripts are available on GitHub under an open-source licence; a Zenodo snapshot will be deposited at publication. Statistics are reported with effect sizes, permutation $p$-values, bootstrap intervals, and Bonferroni adjustment.

\noindent\textbf{Suggested reviewers.}
\begin{enumerate}
  \item Dr.\ Fidele Ntie-Kang, University of Buea --- African natural-product computational discovery
  \item Dr.\ Kelly Chibale, University of Cape Town --- antimalarial drug discovery
  \item Dr.\ John D.\ Chodera, Memorial Sloan Kettering --- free-energy methods and MD best practices
\end{enumerate}

\noindent This manuscript is original, not under consideration elsewhere, and all authors have approved the submission.

\vspace{1em}

\noindent Sincerely,

\vspace{1em}

\noindent \textbf{Myke-Vital Temgoua Sao} \\
(on behalf of all co-authors)

\end{document}
